\section{Discussion}

\paragraph{What the mechanism profile says about Bitcoin.} Bitcoin's proposal discussions are dominated by arguments about the ecosystem, about deployment, about attacks and about incentives. This is the profile of a community whose changes are adopted by third parties rather than shipped by the project: a change must be safe for the nodes that will run it, affordable for the miners and users who will pay for it, and acceptable to the wallets and exchanges that will implement it, and every one of those constituencies can refuse. Authority is rarely invoked because there is little to invoke; the BIP editors administer numbers, the maintainers merge code, and the lead maintainers---while they existed---did not argue. The mechanism the Python paper called \emph{functional}, the argument about whether a proposal is a good idea, is where Python spends a quarter of its influence and Bitcoin a thirteenth: Bitcoin's proposals are less often bad ideas than dangerous or expensive ones.

\paragraph{Governance without a governor.} The era comparison shows what governance looks like when it is contested rather than delegated. The block-size war has the highest compatibility share of any period and a fall in controversial mechanisms rather than a rise; the community argued the war as a compatibility question and, on this measure, argued it more civilly than it argued the early BIPs or the post-lead debates. The post-lead era shows a different signature: security and economic arguments at their peak, the formal roles almost silent, core developers writing a third of all candidates, and controversial mechanisms rising again---driven, as Section~\ref{sec:rq6} shows, by incivility and procedural control around the covenant proposals, the relay-policy disputes and the editorship itself. In Python the transition from a person to a council moved contested influence from fiat to process; in Bitcoin the disappearance of the lead maintainer moved it from a small formal core to a larger informal one, and the process questions that Python's council absorbed are fought on the list.

\paragraph{The decision is not in the repository.} The most consequential methodological result is the failure of the decision-window analysis to transfer. Influence Miner's temporal features assume that a proposal has a decision date; Python's PEP repository supplies one because the BDFL's pronouncement and the status change coincide. Bitcoin's repository supplies a deployment date, often years later and sometimes in bulk. A tool that wants to find the decision in Bitcoin must look elsewhere: at the merge of the implementation, the release that ships it, the start and end of the signalling period, the flag day. The state history built here from the \texttt{bitcoin/bips} repository is the wrong clock for that question, and future work should assemble the right one from the reference implementation's release history and the activation records on the chain.

\paragraph{Volume is not influence, again.} Stance alignment sits at the base rate in Bitcoin as in Python, and the blocking figure is lower still because Bitcoin rarely rejects a BIP formally---contested proposals stay drafts. A decision-mining tool that counted objections would conclude that objections never work; the record of BIP 101, BIP 148 and BIP 119 says otherwise. The mechanisms and roles attached to each sentence are what let an analyst see that the objections that stopped the block-size increase were compatibility and ecosystem arguments from developers, and that the objections that stalled the covenant proposals were security and economic arguments from the community.

\paragraph{Corpus building as a contribution.} The Bitcoin corpus took one working day to build from three public sources, and the same scripts will rebuild it next year. Two of its properties deserve emphasis for other researchers: the public-inbox mirror is the only archive that spans all three homes of the list, and the Google Groups era is unusable without recovering the original sender from \texttt{X-Original-From}. The BIP repository's 2025 relabelling, which collapsed rejected, withdrawn and obsolete proposals into ``Closed'', is a warning that a proposal repository's current state can erase the distinctions a study needs; the git history keeps them.

\section{Data-Quality Findings in the Bitcoin Corpus}
\label{sec:dataquality}

Four properties of the sources would have distorted the results if left untreated. \emph{Release announcements}: the lead maintainer's release notes list every BIP shipped in a release; the announcement of Bitcoin Core 0.21.0 alone produced 1,557 candidate sentences on a first run, and van der Laan appeared as the most influential participant in the corpus. A subject filter now skips the announcements themselves (229 messages) while keeping the replies. \emph{DMARC rewriting}: Google Groups replaces the sender address of every post with the group's own; 687 messages would have been credited to one ``sender'' called \emph{conduition}, the display name of the first such post the loader saw. \emph{Bulk status changes}: 143 of the 505 status changes in the BIP repository happened on 12 April 2025, and 48 more on four other days, so a decision date read naively from the history is often the date of a clean-up. \emph{Technical homonyms}: Bitcoin's vocabulary collides with the controversial cue lists in ways Python's does not---\emph{unilateral} (Lightning's unilateral close), \emph{fork off} and \emph{split the chain} (protocol behaviour), \emph{offline} and \emph{out of band} (wallets and signing), \emph{funds paid to} (outputs), \emph{dumb} (dumb clients). Four passes over the full corpus were needed to add the disambiguating contexts, and the residual inflation of \emph{backchannel} and \emph{ecosystem} by technical uses is noted where it matters.

\section{Threats to Validity}
\label{sec:threats}

\paragraph{Construct validity.} The labels are cue-based and unvalidated; the sample tables show that the cues catch what they were written for and also that they catch technical homonyms. The base mechanisms most affected are \emph{ecosystem} (nodes, wallets and miners as components) and \emph{security} (attack as a protocol term); the rank order of mechanisms is robust to this, the absolute shares are not. The stratified annotation sample that the tool exports is the instrument for the two-coder study.

\paragraph{Coverage.} Only messages that name a BIP number are linked. Proposals are discussed by name before they have a number (``Taproot'', ``CTV'', ``the block size'') and many threads on bitcoin-dev never mention a BIP at all; 6,886 of 24,391 messages are linked. The pre-number phase of a proposal is therefore under-represented, which contributes to the small early-discussion share in RQ1. Title matching, as used in the earlier DeMaP Miner study of BIPs, would raise coverage and is the first extension to make.

\paragraph{Roles.} Rosters were derived from the repositories: maintainers from merge commits (which since 2021 are made by a bot on maintainers' behalf, so recent terms were set open-ended by hand), core developers from a commit threshold, editors and lead maintainers from documented terms. People are matched by clustered display name; pseudonymous participants who changed handles are counted twice. Companies are not resolved: the corporate-interest mechanism sees what people say about employers, not who employs them.

\paragraph{External validity.} The corpus is one mailing list. Bitcoin's decisions are also made on GitHub pull requests, on IRC (whose logs the earlier study used) and, since 2024, on Delving Bitcoin; the list carries the proposal announcements and the set-piece debates, not the code review. Findings are about how influence is exercised where BIPs are discussed publicly.

\section{Conclusion}

This paper applied Influence Miner to the complete public record of Bitcoin's proposal discussions, after rebuilding the corpus from the public-inbox mirror of bitcoin-dev, the BIP repository's history and the reference implementation's commit history. Across 199 BIPs and 62,789 candidate sentences, Bitcoin argues about the ecosystem, deployment, security and incentives where Python argues about implementation and design; its formal leaders barely argue at all; its decisions are recorded by the repository long after they are taken; its most contested period was argued as a compatibility question; and its controversial influence has shifted, since the lead maintainer left, from the formal core to a larger informal one and from fiat to process and incivility. Measured on the same instrument, the two projects differ in what they argue about and who argues, and agree in how explicitly they argue and how little stance volume tells about outcomes. The next steps are the annotation study, title-based linking of pre-number discussions, a decision clock built from releases and activations rather than repository status, the Preference Miner run that RQ5 needs, and the extension of the corpus to the pull-request and Delving Bitcoin venues.

